% ════════════════════════════════════════════════════════════════════════════
% GUÍA — OPERATORIA ENTERA Y FRACCIONARIA · AQUILES · SEMANA 17
% ────────────────────────────────────────────────────────────────────────────
% MATERIAL DEL TUTOR.
% Confirmada tras diagnóstico 2026-04-16: brecha principal en fracciones/%.
% Operatoria entera regular (refuerzo breve). Álgebra y ecuaciones sólidas
% (no se tocan esta semana).
%
% Estructura: ejercicios progresivos + cajas guía post-error. Primeros 10 min
% de la sesión son devolución del diagnóstico. El foco es BLOQUE 2 y 3.
% ════════════════════════════════════════════════════════════════════════════

\documentclass[11pt,a4paper]{article}
\usepackage[utf8]{inputenc}
\usepackage[T1]{fontenc}
\usepackage[spanish]{babel}
\usepackage{geometry}
\usepackage{amsmath,amssymb}
\usepackage{xcolor}
\usepackage{tcolorbox}
\usepackage{enumitem}
\usepackage{fancyhdr}
\usepackage{titlesec}
\usepackage{multirow}
\usepackage{array}
\usepackage{booktabs}
\usepackage{tikz}
\usepackage{hyperref}

\geometry{margin=2cm, top=2.5cm, bottom=2.5cm}

% ─── Colores ───
\definecolor{azulprincipal}{HTML}{2980B9}
\definecolor{azuloscuro}{HTML}{1A5276}
\definecolor{verdequimica}{HTML}{1ABC9C}
\definecolor{verdeclaro}{HTML}{D5F5E3}
\definecolor{naranja}{HTML}{E67E22}
\definecolor{rojo}{HTML}{E74C3C}
\definecolor{amarillorecuerda}{HTML}{FFF3CD}
\definecolor{amarilloborde}{HTML}{FFC107}
\definecolor{grisfondo}{HTML}{F5F5F5}
\definecolor{tutormorado}{HTML}{8E44AD}
\definecolor{tutorfondo}{HTML}{F4ECF7}

% ─── Cajas ───
\tcbset{
  recuerda/.style={
    colback=amarillorecuerda, colframe=amarilloborde,
    fonttitle=\bfseries, title={\color{black} RECUERDA},
    boxrule=1pt, arc=2mm, left=5pt, right=5pt
  },
  propiedad/.style={
    colback=verdeclaro!50, colframe=verdequimica!80!black,
    fonttitle=\bfseries, boxrule=0.8pt, arc=2mm, left=5pt, right=5pt
  },
  formula/.style={
    colback=grisfondo, colframe=azulprincipal,
    boxrule=0.8pt, arc=2mm, left=5pt, right=5pt
  },
  definicion/.style={
    colback=blue!5, colframe=azulprincipal,
    fonttitle=\bfseries, boxrule=0.8pt, arc=2mm, left=5pt, right=5pt
  },
  ejemplo/.style={
    colback=orange!5, colframe=naranja,
    fonttitle=\bfseries, boxrule=0.8pt, arc=2mm, left=5pt, right=5pt
  },
  guia-tutor/.style={
    colback=tutorfondo, colframe=tutormorado,
    fonttitle=\bfseries, boxrule=1pt, arc=2mm, left=5pt, right=5pt,
    title={\color{tutormorado} GUÍA PARA EL TUTOR}
  }
}

% ─── Encabezados ───
\pagestyle{fancy}
\fancyhf{}
\fancyhead[C]{\small\color{gray} Operatoria entera y fraccionaria --- Aquiles --- 2° Medio}
\fancyfoot[C]{\small\color{gray} Página \thepage\ --- \textbf{MATERIAL DEL TUTOR}}
\renewcommand{\headrulewidth}{0.4pt}

% ─── Títulos ───
\newcommand{\tituloseccion}[1]{%
  \noindent\colorbox{azulprincipal}{%
    \makebox[\dimexpr\textwidth-2\fboxsep\relax][l]{%
      \color{white}\normalfont\Large\bfseries\ BLOQUE \thesection: #1%
    }%
  }%
}
\titleformat{\section}
  {}{}{0em}
  {\tituloseccion}
\titleformat{\subsection}
  {\normalfont\large\bfseries\color{azulprincipal}}
  {\thesubsection}{0.5em}{}

% ─── Niveles ───
\newcommand{\nivelbasico}{%
  \par\smallskip\colorbox{green!60!black}{\color{white}\small\bfseries\ NIVEL BÁSICO\ }\par\smallskip}
\newcommand{\nivelintermedio}{%
  \par\smallskip\colorbox{naranja}{\color{white}\small\bfseries\ NIVEL INTERMEDIO\ }\par\smallskip}

% ─── Observación ───
\newcommand{\observacion}[1]{%
  \par\medskip
  \noindent\colorbox{grisfondo}{%
    \makebox[\dimexpr\textwidth-2\fboxsep\relax][l]{%
      \small\color{gray}\itshape\ #1%
    }%
  }%
  \par\medskip
}

% ─── Líneas ───
\newcommand{\lineasrespuesta}[1]{%
  \par\vspace{2mm}%
  \foreach \i in {1,...,#1}{\noindent\rule{\textwidth}{0.3pt}\par\vspace{5mm}}%
  \vspace{2mm}%
}

\begin{document}

% ════════════════════════
% PORTADA
% ════════════════════════
\thispagestyle{empty}
\vspace*{2.5cm}
\begin{center}
  {\Huge\bfseries\color{azulprincipal} GUÍA DE NIVELACIÓN}\\[8pt]
  {\LARGE\bfseries\color{azuloscuro} Operatoria entera y fraccionaria}\\[15pt]
  {\color{azulprincipal}\rule{8cm}{2pt}}\\[15pt]
  {\Large Semana 17 \textbar{} 2° Medio \textbar{} 2026-S1}\\[8pt]
  {\large Alumno: Aquiles --- Sesión 2026-04-23}\\[25pt]
  {\itshape\color{gray}
  Material del tutor. Confirmado tras diagnóstico 2026-04-16.\\[4pt]
  Brecha principal: fracciones y \%. Operatoria: refuerzo breve.\\[4pt]
  Álgebra y ecuaciones sólidas --- no se tocan esta semana.}
\end{center}

\vspace{1cm}
\begin{tcolorbox}[recuerda]
\textbf{Plan de la sesión (60 min):}
\begin{enumerate}[nosep]
  \item \textbf{Devolución del diagnóstico (10 min):} nombrar lo que hizo bien (álgebra, ecuaciones), ubicar la brecha (fracciones) sin dramatizar, mostrar el plan de hoy.
  \item \textbf{Bloque A --- Signos y jerarquía (10 min):} Bloque 1, ejercicios 1--5 y 11--13. Afianzar, no re-enseñar. Énfasis en $(-4)^2$ vs. $-4^2$.
  \item \textbf{Bloque B --- Fracciones (25 min):} Bloque 2, partir del dibujo. Simplificación $\to$ suma $\to$ multiplicación. \textbf{Foco principal.}
  \item \textbf{Bloque C --- Porcentajes directos (10 min):} Bloque 3, tabla de equivalencias + atajos mentales + camino A de descuento.
  \item \textbf{Cierre + tarea (5 min):} tarea corta, preguntar cómo se sintió, ajustar plan de la semana 18.
\end{enumerate}
\end{tcolorbox}

\newpage

% ════════════════════════════════════════════════════
% BLOQUE 1 — ENTEROS CON SIGNO
% ════════════════════════════════════════════════════
\section{ENTEROS CON SIGNO}

\observacion{Usar este bloque si el diagnóstico mostró brecha en la regla de signos o en sumas/restas con negativos.}

\subsection{Ejercicios progresivos}

\nivelbasico

\textbf{Sumas y restas} (dictar o escribir en pizarra):
\begin{enumerate}[label=\textbf{\arabic*.},leftmargin=1cm,itemsep=6pt]
  \item $(-3) + 8 = $ \hfill \textit{($5$)}
  \item $(-5) + (-4) = $ \hfill \textit{($-9$)}
  \item $7 - 12 = $ \hfill \textit{($-5$)}
  \item $(-6) - (-2) = $ \hfill \textit{($-4$)}
  \item $(-1) + 9 - 3 - (-5) = $ \hfill \textit{($10$)}
\end{enumerate}

\medskip
\textbf{Multiplicación y división:}
\begin{enumerate}[label=\textbf{\arabic*.},leftmargin=1cm,itemsep=6pt,start=6]
  \item $(-3) \cdot 7 = $ \hfill \textit{($-21$)}
  \item $(-8) \cdot (-4) = $ \hfill \textit{($32$)}
  \item $(-36) \div 9 = $ \hfill \textit{($-4$)}
  \item $(-45) \div (-5) = $ \hfill \textit{($9$)}
  \item $(-2) \cdot (-3) \cdot (-1) = $ \hfill \textit{($-6$)}
\end{enumerate}

\nivelintermedio

\textbf{Operaciones combinadas} (respetar jerarquía):
\begin{enumerate}[label=\textbf{\arabic*.},leftmargin=1cm,itemsep=6pt,start=11]
  \item $3 + (-2) \cdot 5 = $ \hfill \textit{($-7$)}
  \item $(-4)^{2} - 10 = $ \hfill \textit{($6$)}
  \item $-4^{2} - 10 = $ \hfill \textit{($-26$)  --- ¡ojo con la diferencia!}
  \item $12 \div (-3) + (-2) \cdot 4 = $ \hfill \textit{($-12$)}
  \item $(-1)^{5} + (-1)^{4} = $ \hfill \textit{($0$)}
\end{enumerate}

\subsection{Guías post-error}

\begin{tcolorbox}[guia-tutor, title={\color{tutormorado} RESTAR UN NEGATIVO}]
``$(-6) - (-2)$'' se lee ``menos seis \textbf{menos} menos dos''. Los dos signos negativos seguidos se transforman en un más:
\[
  (-6) - (-2) = (-6) + 2 = -4
\]
\textbf{Truco mnemotécnico:} dos signos iguales juntos $\to$ positivo. Dos distintos $\to$ negativo.

Practicar con 3 ejemplos más: $10 - (-3)$, $(-7) - (-7)$, $(-1) - (-8)$.
\end{tcolorbox}

\begin{tcolorbox}[guia-tutor, title={\color{tutormorado} DIFERENCIA ENTRE $(-4)^{2}$ y $-4^{2}$}]
Este es uno de los errores más comunes de todo el currículo. Escribirlo en grande:

\begin{center}
\renewcommand{\arraystretch}{1.5}
\begin{tabular}{@{}ll@{}}
$(-4)^{2} = (-4)\cdot(-4) = 16$ & El paréntesis incluye el signo en la base. \\
$-4^{2} = -(4 \cdot 4) = -16$ & El signo no está en la base; el cuadrado se aplica solo al $4$. \\
\end{tabular}
\end{center}

Si no lo ve: ``¿qué está elevado al cuadrado? ¿El $-4$ completo o solo el $4$?'' El paréntesis decide.
\end{tcolorbox}

\begin{tcolorbox}[guia-tutor, title={\color{tutormorado} PRODUCTO CON CANTIDAD IMPAR DE NEGATIVOS}]
En $(-2) \cdot (-3) \cdot (-1)$:
\begin{enumerate}[nosep]
  \item Primero: $(-2) \cdot (-3) = 6$ (dos negativos $\to$ positivo).
  \item Luego: $6 \cdot (-1) = -6$ (positivo $\times$ negativo $\to$ negativo).
\end{enumerate}

\textbf{Regla rápida:} contar los factores negativos. Si son \textbf{par} $\to$ resultado positivo. Si son \textbf{impar} $\to$ negativo. No necesita multiplicar todo para saber el signo.
\end{tcolorbox}

\newpage

% ════════════════════════════════════════════════════
% BLOQUE 2 — FRACCIONES
% ════════════════════════════════════════════════════
\section{FRACCIONES}

\observacion{Usar este bloque si el diagnóstico mostró que las fracciones son ``dos números apilados'' y no una cantidad con significado.}

\subsection{Ejercicios progresivos}

\nivelbasico

\textbf{Simplificar:}
\begin{enumerate}[label=\textbf{\arabic*.},leftmargin=1cm,itemsep=6pt]
  \item $\dfrac{12}{18} = $ \hfill \textit{$\left(\dfrac{2}{3}\right)$}
  \item $\dfrac{20}{35} = $ \hfill \textit{$\left(\dfrac{4}{7}\right)$}
  \item $\dfrac{24}{36} = $ \hfill \textit{$\left(\dfrac{2}{3}\right)$}
\end{enumerate}

\medskip
\textbf{Sumar y restar:}
\begin{enumerate}[label=\textbf{\arabic*.},leftmargin=1cm,itemsep=6pt,start=4]
  \item $\dfrac{2}{5} + \dfrac{1}{5} = $ \hfill \textit{$\left(\dfrac{3}{5}\right)$}
  \item $\dfrac{1}{2} + \dfrac{1}{3} = $ \hfill \textit{$\left(\dfrac{5}{6}\right)$}
  \item $\dfrac{3}{4} - \dfrac{1}{6} = $ \hfill \textit{$\left(\dfrac{7}{12}\right)$}
  \item $1 - \dfrac{2}{5} = $ \hfill \textit{$\left(\dfrac{3}{5}\right)$}
\end{enumerate}

\medskip
\textbf{Multiplicar y dividir:}
\begin{enumerate}[label=\textbf{\arabic*.},leftmargin=1cm,itemsep=6pt,start=8]
  \item $\dfrac{2}{3} \cdot \dfrac{5}{4} = $ \hfill \textit{$\left(\dfrac{5}{6}\right)$}
  \item $\dfrac{3}{7} \div \dfrac{6}{5} = $ \hfill \textit{$\left(\dfrac{5}{14}\right)$}
\end{enumerate}

\nivelintermedio

\textbf{Operaciones combinadas:}
\begin{enumerate}[label=\textbf{\arabic*.},leftmargin=1cm,itemsep=6pt,start=10]
  \item $\dfrac{1}{2} + \dfrac{1}{3} \cdot \dfrac{3}{4} = $ \hfill \textit{$\left(\dfrac{3}{4}\right)$} {\small --- ¡jerarquía!}
  \item $\dfrac{2}{3} \cdot \left(\dfrac{1}{2} + \dfrac{1}{4}\right) = $ \hfill \textit{$\left(\dfrac{1}{2}\right)$}
\end{enumerate}

\subsection{Guías post-error}

\begin{tcolorbox}[guia-tutor, title={\color{tutormorado} SI NO ``VE'' QUÉ ES UNA FRACCIÓN}]
Dibujar un círculo (o una barra) y dividirlo:

\begin{center}
\begin{tikzpicture}[scale=0.7]
  \draw[thick] (0,0) circle (1.5cm);
  \draw[thick] (0,0) -- (0,1.5);
  \draw[thick] (0,0) -- (1.3,-0.75);
  \draw[thick] (0,0) -- (-1.3,-0.75);
  \fill[azulprincipal!30] (0,0) -- (0,1.5) arc (90:210:1.5) -- cycle;
  \node at (3,0) {$= \dfrac{2}{3}$ de la pizza};
\end{tikzpicture}
\end{center}

``$\dfrac{2}{3}$'' significa: dividí en \textbf{3 partes iguales} y tomé \textbf{2}. El de abajo dice \emph{en cuántas partes corté}; el de arriba dice \emph{cuántas me quedé}.

Que él dibuje $\dfrac{3}{4}$ y $\dfrac{1}{5}$ para verificar que lo entiende.
\end{tcolorbox}

\begin{tcolorbox}[guia-tutor, title={\color{tutormorado} SI SUMA NUMERADORES \emph{Y} DENOMINADORES}]
El error $\dfrac{1}{2} + \dfrac{1}{3} = \dfrac{2}{5}$ es el más frecuente de toda la secundaria.

\textbf{Contraejemplo visual:} ``Si me como media pizza y luego un tercio de pizza, ¿me comí dos quintos de pizza? No --- me comí \emph{más} que media. Dos quintos es \emph{menos} que media. Algo está mal.''

Procedimiento correcto:
\[
  \frac{1}{2} + \frac{1}{3} = \frac{1 \times 3}{2 \times 3} + \frac{1 \times 2}{3 \times 2} = \frac{3}{6} + \frac{2}{6} = \frac{5}{6}
\]
Lo que se amplifica es la fracción completa (arriba \textbf{y} abajo por el mismo número). El denominador no se suma: solo se iguala.
\end{tcolorbox}

\begin{tcolorbox}[guia-tutor, title={\color{tutormorado} SI NO SABE DIVIDIR FRACCIONES}]
\textbf{Dividir fracciones = multiplicar por la inversa} (dar vuelta la segunda).
\[
  \frac{3}{7} \div \frac{6}{5} = \frac{3}{7} \cdot \frac{5}{6} = \frac{15}{42} = \frac{5}{14}
\]

Mnemotecnia: ``\textbf{dividir es dar vuelta y multiplicar}''. Verificar con un ejemplo entero primero: $6 \div 2 = 6 \cdot \frac{1}{2} = 3$.
\end{tcolorbox}

\begin{tcolorbox}[guia-tutor, title={\color{tutormorado} SI NO RESPETA JERARQUÍA CON FRACCIONES}]
En $\dfrac{1}{2} + \dfrac{1}{3} \cdot \dfrac{3}{4}$:

La multiplicación se hace \textbf{primero}:
\[
  \frac{1}{3} \cdot \frac{3}{4} = \frac{3}{12} = \frac{1}{4}
\]
Luego la suma:
\[
  \frac{1}{2} + \frac{1}{4} = \frac{2}{4} + \frac{1}{4} = \frac{3}{4}
\]

Si no lo ve: ``la multiplicación es más fuerte que la suma --- siempre va primero, igual que con enteros''.
\end{tcolorbox}

\newpage

% ════════════════════════════════════════════════════
% BLOQUE 3 — CONEXIÓN FRACCIÓN ↔ DECIMAL ↔ %
% ════════════════════════════════════════════════════
\section{EQUIVALENCIAS: FRACCIÓN, DECIMAL Y PORCENTAJE}

\observacion{Este bloque complementa el anterior. Si las fracciones están sólidas, se puede pasar directo a la tabla de equivalencias y los ejercicios de \%.}

\subsection{Tabla de referencia}

\begin{tcolorbox}[formula]
\begin{center}
\renewcommand{\arraystretch}{1.4}
\begin{tabular}{@{}ccc@{}}
\toprule
\textbf{Fracción} & \textbf{Decimal} & \textbf{Porcentaje} \\
\midrule
$\dfrac{1}{2}$ & $0{,}5$ & $50\%$ \\[4pt]
$\dfrac{1}{4}$ & $0{,}25$ & $25\%$ \\[4pt]
$\dfrac{3}{4}$ & $0{,}75$ & $75\%$ \\[4pt]
$\dfrac{1}{5}$ & $0{,}2$ & $20\%$ \\[4pt]
$\dfrac{1}{10}$ & $0{,}1$ & $10\%$ \\[4pt]
$\dfrac{1}{3}$ & $0{,}\overline{3}$ & $33{,}\overline{3}\%$ \\
\bottomrule
\end{tabular}
\end{center}
\end{tcolorbox}

\subsection{Ejercicios}

\textbf{Completar la equivalencia:}
\begin{enumerate}[label=\textbf{\arabic*.},leftmargin=1cm,itemsep=6pt]
  \item $\dfrac{3}{5} = $ \rule{1.5cm}{0.4pt} $= $ \rule{1.5cm}{0.4pt}$\%$ \hfill \textit{($0{,}6$ / $60\%$)}
  \item $0{,}4 = \dfrac{\qquad}{\qquad} = $ \rule{1.5cm}{0.4pt}$\%$ \hfill \textit{$\left(\dfrac{2}{5}\right.$ / $\left.40\%\right)$}
  \item $75\% = \dfrac{\qquad}{\qquad} = $ \rule{1.5cm}{0.4pt} \hfill \textit{$\left(\dfrac{3}{4}\right.$ / $\left.0{,}75\right)$}
\end{enumerate}

\medskip
\textbf{Problemas de \%:}
\begin{enumerate}[label=\textbf{\arabic*.},leftmargin=1cm,itemsep=6pt,start=4]
  \item Un pantalón cuesta $\$30\,000$ y tiene un $20\%$ de descuento. ¿Cuánto se paga? \hfill \textit{($\$24\,000$)}
  \item En una prueba de 40 preguntas, acertaste el $70\%$. ¿Cuántas acertaste? \hfill \textit{($28$)}
  \item Un producto subió de $\$15\,000$ a $\$18\,000$. ¿Qué porcentaje subió? \hfill \textit{($20\%$)}
\end{enumerate}

\subsection{Guías post-error}

\begin{tcolorbox}[guia-tutor, title={\color{tutormorado} CÓMO CONVERTIR ENTRE LAS TRES FORMAS}]
\textbf{Fracción $\to$ decimal:} dividir numerador entre denominador. $\frac{3}{5} = 3 \div 5 = 0{,}6$.

\textbf{Decimal $\to$ porcentaje:} multiplicar por 100. $0{,}6 \times 100 = 60\%$.

\textbf{Porcentaje $\to$ fracción:} poner sobre 100 y simplificar. $75\% = \frac{75}{100} = \frac{3}{4}$.

El ciclo completo:
\[
  \text{fracción} \xrightarrow{\div} \text{decimal} \xrightarrow{\times 100} \text{porcentaje} \xrightarrow{\div 100} \text{fracción}
\]
\end{tcolorbox}

\begin{tcolorbox}[guia-tutor, title={\color{tutormorado} SI NO PUEDE CALCULAR UN \% DE AUMENTO/DESCUENTO}]
Dos caminos (que él elija el que le sea más cómodo):

\textbf{Camino A} (calcular el \% y sumar/restar):
\begin{itemize}[nosep]
  \item $20\%$ de $30\,000 = 0{,}2 \times 30\,000 = 6\,000$.
  \item Descuento: $30\,000 - 6\,000 = 24\,000$.
\end{itemize}

\textbf{Camino B} (multiplicar por el factor directo):
\begin{itemize}[nosep]
  \item $20\%$ de descuento $\to$ queda el $80\%$ $\to$ factor $0{,}8$.
  \item $30\,000 \times 0{,}8 = 24\,000$.
\end{itemize}

El camino B es más eficiente, pero el A es más intuitivo. En esta etapa, que use el A. El B puede venir después.
\end{tcolorbox}

\begin{tcolorbox}[guia-tutor, title={\color{tutormorado} SI NO PUEDE CALCULAR QUÉ \% REPRESENTA UN CAMBIO}]
De $\$15\,000$ a $\$18\,000$:
\begin{enumerate}[nosep]
  \item ¿Cuánto subió? $18\,000 - 15\,000 = 3\,000$.
  \item ¿Ese $3\,000$ qué fracción es del original? $\dfrac{3\,000}{15\,000} = \dfrac{1}{5} = 0{,}2$.
  \item Convertir a \%: $0{,}2 \times 100 = 20\%$.
\end{enumerate}

\textbf{Clave:} el porcentaje de cambio es \textbf{siempre respecto al valor original} (el de antes), no respecto al nuevo.
\end{tcolorbox}

\newpage

% ════════════════════════════════════════════════════
% HOJA DE PROGRESO
% ════════════════════════════════════════════════════
\section*{NOTAS DE PROGRESO --- SEMANA 17}
\addcontentsline{toc}{section}{Notas de progreso}

\begin{center}
{\large\bfseries\color{azuloscuro} Aquiles --- Sesión del \rule{3cm}{0.4pt} de abril 2026}
\end{center}

\medskip

\noindent\textbf{Devolución del diagnóstico --- ¿cómo lo tomó?}
\lineasrespuesta{3}

\medskip

\noindent\textbf{Bloque trabajado hoy:} \rule{5cm}{0.4pt}

\medskip

\noindent\textbf{Lo que hizo bien:}
\lineasrespuesta{2}

\noindent\textbf{Lo que costó:}
\lineasrespuesta{2}

\noindent\textbf{Nivel emocional durante la sesión} (frustración / enganche / indiferencia):
\lineasrespuesta{2}

\noindent\textbf{Tarea asignada:}
\lineasrespuesta{2}

\noindent\textbf{¿Se confirma el tema de la semana 18 o hay que ajustar?}
\lineasrespuesta{2}

\begin{tcolorbox}[recuerda]
Trasladar estas notas a la sección ``Notas para la próxima semana'' del plan semanal correspondiente en el vault.
\end{tcolorbox}

\end{document}
